\section{Finite matching lower-bound statements}
\label{app:results}

This section defines the hard matching family and states its finite TV
query--round lower bound and target-error tradeoff.
Appendix~\ref{app:lower-proof} proves these results and verifies class
membership in \cref{app:prop:hard-family-inclusion}.
Appendix~\ref{app:benchmark-specializations} then places the finite
main-text calibration (\cref{app:cor:signal-calibrated-finite-lower})
immediately before the fixed-accuracy derivation
(\cref{app:cor:fixed-accuracy-main}) of the \HellingerCase\ lower bound
in \cref{main:thm:lower-envelope}.
The finite upper bounds for the two cases are
\cref{app:thm:finite-upper,app:thm:finite-kl-upper}.

\subsection{Finite matching lower bound}
\label{app:results-lower}

\begin{definition}[Finite hard matching family]
\label{app:def:hard-matching-family}
Let $N$ be even, let
$1\le\HardBankSize<\HardVocabSize$ be integers, and suppose
\begin{equation}
 \HardBankSize\TriggerMass<1,
 \qquad
 0<\EdgeCoupling\le\TriggerMass\le\frac12.
 \label{app:eq:hard-basic-conditions}
\end{equation}
Use the public vocabulary and common marginal
\begin{equation}
 \Vocab=\{\HardCandidate{1},\ldots,\HardCandidate{\HardBankSize}\}
        \sqcup\HardBackground,
 \qquad |\HardBackground|=\HardVocabSize-\HardBankSize,
 \qquad
 \HardMarginal(\HardCandidate{c})=\TriggerMass,
 \qquad
 \HardMarginal(b)=
 \frac{1-\HardBankSize\TriggerMass}
      {\HardVocabSize-\HardBankSize}.
 \label{app:eq:hard-common-marginal}
\end{equation}
The candidate bank
$\{\HardCandidate{1},\ldots,\HardCandidate{\HardBankSize}\}$ and the
background set are public.  An instance
$I=(\HiddenMatching,\TriggerVector)$ consists of a hidden perfect matching
$\HiddenMatching$ on $\PositionSet$ and hidden trigger indices
$\TriggerVector=(\TriggerIndex{i})_{i=1}^N
\in[\HardBankSize]^N$.  Put
\begin{equation}
 \TriggerFeature{i}(a)
 :=\mathbf 1\{a=\HardCandidate{\TriggerIndex{i}}\}-\TriggerMass,
 \qquad
 \TriggerVariance:=\TriggerMass(1-\TriggerMass).
 \label{app:eq:hard-features}
\end{equation}
For $e=\{i,j\}\in\HiddenMatching$, define
\begin{align}
 \HardIndependentEdgeLaw{e}(a,b)
 &:=\HardMarginal(a)\HardMarginal(b),
 \label{app:eq:hard-independent-edge-law}\\
 \HardEdgeLaw{e}{I}(a,b)
 &:=\HardMarginal(a)\HardMarginal(b)
 \left[1+\EdgeCoupling
 \frac{\TriggerFeature{i}(a)\TriggerFeature{j}(b)}{\TriggerVariance}\right],
 \label{app:eq:hard-edge-law}\\
 \HardTargetLaw{I}&:=\bigotimes_{e\in\HiddenMatching}\HardEdgeLaw{e}{I},
 \qquad
 \HardKernel{i}{j}{I}(b\mid a)
 :=\HardMarginal(b)
 \left[1+\EdgeCoupling
 \frac{\TriggerFeature{i}(a)\TriggerFeature{j}(b)}{\TriggerVariance}\right].
 \label{app:eq:hard-target-and-kernel}
\end{align}
The bounds in \eqref{app:eq:hard-basic-conditions} make these laws strictly
positive: the three possible likelihood-ratio values are displayed in
\eqref{app:eq:hard-likelihood-values} and checked in Step~1 of the proof of
Lemma~\ref{app:lem:hard-local-facts}.
If $\MateInMatching{j}$ denotes the mate of $j$, the selected
frozen oracle is, for every masked readout $j$ and
$i:=\MateInMatching{j}$,
\begin{equation}
 \HardOracleRow{I}{y}{j}
 :=\begin{cases}
 \HardMarginal,
   &i\in\MaskSet{y},\\
 \HardKernel{i}{j}{I}(\cdot\mid y_i),
   &i\in\ObservedSet{y}\ \text{and}\
     y_i=\HardCandidate{\TriggerIndex{i}},\\
 \HardMarginal,
   &i\in\ObservedSet{y}\ \text{and}\
     y_i\ne\HardCandidate{\TriggerIndex{i}}.
 \end{cases}
 \label{app:eq:hard-oracle}
\end{equation}
Let
$\HardFamily[\HardVocabSize,\HardBankSize,\TriggerMass,\EdgeCoupling]$ be the
family obtained by ranging over all $\HiddenMatching$ and $\TriggerVector$.
The public candidate labels are granted to the algorithm;
the matching and trigger indices remain hidden.
\end{definition}

For integers $\QueryBudget{},\RoundBudget{}$, let $\HardAlgorithmClass{\QueryBudget{}}{\RoundBudget{}}$ be the admissible algorithms
whose counterfactual-submission and nonempty-commit counters are at most $\QueryBudget{}$
and $\RoundBudget{}$, respectively, on every instance, seed, transcript, and realized
path.  Define
\begin{equation}
 \HardMinimax(\QueryBudget{},\RoundBudget{};\HardVocabSize,\HardBankSize,
              \TriggerMass,\EdgeCoupling)
 :=\inf_{\mathcal A\in\HardAlgorithmClass{\QueryBudget{}}{\RoundBudget{}}}
   \sup_{I\in\HardFamily[\HardVocabSize,\HardBankSize,
                           \TriggerMass,\EdgeCoupling]}
   \SeedRisk(\mathcal A;\HardTargetLaw{I},\HardOracle{I}).
 \label{app:eq:hard-minimax-risk}
\end{equation}

The theorem below uses only the finite hard-family conditions.
\Cref{app:cor:signal-calibrated-finite-lower} later chooses these parameters
to satisfy the \HellingerCase\ oracle condition for the main lower bound.

\begin{theorem}[Nonlinear finite adaptive query--round lower bound]
\label{app:thm:nonlinear-finite-lower}
Assume only the basic hard-family conditions
\eqref{app:eq:hard-basic-conditions}; in particular, no restriction on
\(N\EdgeCoupling^2\) is imposed. Then
\begin{equation}
 \begin{aligned}
 &0\le \QueryBudget{}\le\HardBankSize/8,
 \qquad
 1\le \RoundBudget{}\le N/16384\\
 &\quad\Longrightarrow\quad
 \HardMinimax(\QueryBudget{},\RoundBudget{};\HardVocabSize,\HardBankSize,
              \TriggerMass,\EdgeCoupling)\\
 &\hspace{7em}\ge
 \frac14\left[
  1-\exp\left\{-\frac{N\EdgeCoupling^2}{98304\RoundBudget{}}\right\}
 \right].
 \end{aligned}
 \label{app:eq:nonlinear-finite-lower}
\end{equation}
\end{theorem}

\begin{corollary}[Nonlinear finite target-error tradeoff]
\label{app:cor:nonlinear-finite-target-tradeoff}
Assume \eqref{app:eq:hard-basic-conditions}, fix
\(\varepsilon\in(0,1/8]\), and suppose an admissible algorithm has
worst-case pathwise budgets $(\QueryBudget{},\RoundBudget{})$ and
\begin{equation}
 \sup_{I\in\HardFamily[\HardVocabSize,\HardBankSize,
                         \TriggerMass,\EdgeCoupling]}
 \SeedRisk(\mathcal A;\HardTargetLaw{I},\HardOracle{I})\le\varepsilon.
 \label{app:eq:finite-target-risk}
\end{equation}
Then at least one of
\begin{equation}
 \QueryBudget{}>\HardBankSize/8,
 \qquad
 \RoundBudget{}>N/16384,
 \qquad
 \RoundBudget{}\ge\frac{N\EdgeCoupling^2}{786432\varepsilon}
 \label{app:eq:nonlinear-finite-target-alternative}
\end{equation}
holds.
\end{corollary}

% Please keep the finite main-text calibration beside its fixed-accuracy
% specialization in appendix_benchmarks.tex.

\begin{remark}[Where the oracle slack enters]
\label{app:rem:oracle-slack}
The selected oracle \eqref{app:eq:hard-oracle} is exact except in its third
case, where a revealed nontrigger source value $a$ at the mate $i$ of $j$
produces the row $\HardMarginal$ instead of the exact row
$\HardKernel{i}{j}{I}(\cdot\mid a)$.  The difference of these two rows is
$\HardMarginal(b)\EdgeCoupling\TriggerFeature{i}(a)\TriggerFeature{j}(b)
/\TriggerVariance$ with $\TriggerFeature{i}(a)=-\TriggerMass$, so by
\eqref{app:eq:hard-feature-moments} they are at total-variation distance
exactly $\EdgeCoupling\TriggerMass$, and the squared Hellinger bound
\eqref{app:eq:hard-local-oracle-error} is the only place where
Assumption~\ref{app:ass:oracle-accuracy} is consumed.  The resulting
constraint $\EdgeCoupling^2\TriggerMass\le(\OracleRadius)^2/(2N)$ in
\eqref{app:eq:finite-hard-inclusion} yields the cubic signal calibration
when $\EdgeCoupling=\TriggerMass=\EdgeSignal$.
\end{remark}

% Please keep the finite upper statement at the start of appendix_upper_proofs.tex.

% Please keep hard-family inclusion with its proof in proofs/lower_family.tex.
